% GENERATED FILE. Edit canonical lessons, then run npm run generate:books.
% canonical-source-hash: fnv1a64:08bf727ececc0e99
% canonical-lessons: ES-C496-hogar, ES-C496-cerradura, ES-C496-colocar, ES-C496-quitar, ES-C496-llenar

\chapter{Cosas de casa}
\label{ch:es-cosas-de-casa}
\hlchaptermodality{496}

\begin{chapteropening}
\textbf{By the end of this chapter:} \emph{I can talk about home: the lock, putting things in place, taking things away and filling something up.}

Complete the last lesson of chapter 496: I can talk about home: the lock, putting things in place, taking things away and filling something up.
\end{chapteropening}

\section[el hogar]{el hogar — the home}

\label{lesson:ES-C496-hogar}



\begin{quote}
Before the new one: say the Spanish for to disappear, then the Spanish for the thief.
\end{quote}

\subsection*{You'll want to know: el hogar}
\textbf{el hogar} — \textquotedblleft{}the home\textquotedblright{}.

\begin{quote}
Productos para el hogar.
\end{quote}

\begin{quote}
\emph{Products for the home.}
\end{quote}

A warmer word than \emph{la casa}; shop signs use it.

\begin{cousinweb}
\textbf{hogar} $\leftarrow$ Old Spanish \emph{fogar} (EtymDB 2.1), from spoken Latin \emph{*focaris}, built on \emph{focus}, \textquotedblleft{}a hearth\textquotedblright{} — inherited: \emph{f-} became \emph{h-}. A home was named for its fire.
\end{cousinweb}

\begin{grammarlens}[title={What you've built}]
One more word for reading notices, forms and messages.
\end{grammarlens}

\subsection*{Your turn}
\begin{itemize}
\raggedright
  \item \emph{Say it:} \emph{el hogar}
  \item \emph{Say it:} \emph{el hogar}, once more
  \item \emph{Say it:} a sentence of your own with \emph{el hogar}
  \item \emph{Recall:} say the Spanish for to disappear, then the Spanish for the thief, then say \emph{el hogar} again
\end{itemize}

\begin{tcolorbox}[breakable,colback=teal!4,colframe=teal!35!black,title={Before you move on}]
What does el hogar mean? (\textquotedblleft{}The home\textquotedblright{}.) Say it once more.
\end{tcolorbox}
\section[la cerradura]{la cerradura — the lock}

\label{lesson:ES-C496-cerradura}



\begin{quote}
Before the new one: say the Spanish for the thief, then the Spanish for the home.
\end{quote}

\subsection*{You'll want to know: la cerradura}
\textbf{la cerradura} — \textquotedblleft{}the lock\textquotedblright{}.

\begin{quote}
La llave no entra en la cerradura.
\end{quote}

\begin{quote}
\emph{The key does not go into the lock.}
\end{quote}

\begin{cousinweb}
\textbf{cerradura} is \emph{cerrar}, \textquotedblleft{}to close\textquotedblright{}, with \emph{-dura}, the ending for the thing that does the job. \emph{Cerrar} is Late Latin \emph{serare}, \textquotedblleft{}to bolt\textquotedblright{}, from \emph{sera}, the bar across a door.
\end{cousinweb}

\begin{grammarlens}[title={What you've built}]
One more word for reading notices, forms and messages.
\end{grammarlens}

\subsection*{Your turn}
\begin{itemize}
\raggedright
  \item \emph{Say it:} \emph{la cerradura}
  \item \emph{Say it:} \emph{la cerradura}, once more
  \item \emph{Say it:} a sentence of your own with \emph{la cerradura}
  \item \emph{Recall:} say the Spanish for the thief, then the Spanish for the home, then say \emph{la cerradura} again
\end{itemize}

\begin{tcolorbox}[breakable,colback=teal!4,colframe=teal!35!black,title={Before you move on}]
What does la cerradura mean? (\textquotedblleft{}The lock\textquotedblright{}.) Say it once more.
\end{tcolorbox}
\section[colocar]{colocar — to put, to place}

\label{lesson:ES-C496-colocar}



\begin{quote}
Before the new one: say the Spanish for the home, then the Spanish for the lock.
\end{quote}

\subsection*{You'll want to know: colocar}
\textbf{colocar} — \textquotedblleft{}to put, to place\textquotedblright{}.

\begin{quote}
Coloca los libros en la estantería.
\end{quote}

\begin{quote}
\emph{Put the books on the shelf.}
\end{quote}

It takes \emph{qu} before \emph{e}: \emph{coloqué}.

\begin{cousinweb}
\textbf{colocar} $\leftarrow$ Latin \emph{collocare}, \textquotedblleft{}to place, to set in order\textquotedblright{} (EtymDB 2.1): \emph{con-} and \emph{locare}, \textquotedblleft{}to put\textquotedblright{} — a learned borrowing. Worn down in speech, the same verb became \emph{colgar}, \textquotedblleft{}to hang\textquotedblright{}.
\end{cousinweb}

\begin{grammarlens}[title={What you've built}]
One more word for reading notices, forms and messages.
\end{grammarlens}

\subsection*{Your turn}
\begin{itemize}
\raggedright
  \item \emph{Say it:} \emph{colocar}
  \item \emph{Say it:} \emph{colocar}, once more
  \item \emph{Say it:} a sentence of your own with \emph{colocar}
  \item \emph{Recall:} say the Spanish for the home, then the Spanish for the lock, then say \emph{colocar} again
\end{itemize}

\begin{tcolorbox}[breakable,colback=teal!4,colframe=teal!35!black,title={Before you move on}]
What does colocar mean? (\textquotedblleft{}To put, to place\textquotedblright{}.) Say it once more.
\end{tcolorbox}
\section[quitar]{quitar — to take away, to remove}

\label{lesson:ES-C496-quitar}



\begin{quote}
Before the new one: say the Spanish for the lock, then the Spanish for to put.
\end{quote}

\subsection*{You'll want to know: quitar}
\textbf{quitar} — \textquotedblleft{}to take away, to remove\textquotedblright{}.

\begin{quote}
Quita la mesa después de comer.
\end{quote}

\begin{quote}
\emph{Clear the table after eating.}
\end{quote}

\emph{Quitarse la chaqueta} — to take off your jacket.

\begin{cousinweb}
\textbf{quitar} $\leftarrow$ Medieval Latin \emph{quitare}, \textquotedblleft{}to set free from a debt\textquotedblright{}, from Latin \emph{quietus}, \textquotedblleft{}at rest\textquotedblright{}. English \emph{quit} is the same word.
\end{cousinweb}

\begin{grammarlens}[title={What you've built}]
One more word for reading notices, forms and messages.
\end{grammarlens}

\subsection*{Your turn}
\begin{itemize}
\raggedright
  \item \emph{Say it:} \emph{quitar}
  \item \emph{Say it:} \emph{quitar}, once more
  \item \emph{Say it:} a sentence of your own with \emph{quitar}
  \item \emph{Recall:} say the Spanish for the lock, then the Spanish for to put, then say \emph{quitar} again
\end{itemize}

\begin{tcolorbox}[breakable,colback=teal!4,colframe=teal!35!black,title={Before you move on}]
What does quitar mean? (\textquotedblleft{}To take away, to remove\textquotedblright{}.) Say it once more.
\end{tcolorbox}
\section[llenar]{llenar — to fill}

\label{lesson:ES-C496-llenar}



\begin{quote}
Before the new one: say the Spanish for to put, then the Spanish for to take away.
\end{quote}

\subsection*{You'll want to know: llenar}
\textbf{llenar} — \textquotedblleft{}to fill\textquotedblright{}.

\begin{quote}
Llena la botella de agua.
\end{quote}

\begin{quote}
\emph{Fill the bottle with water.}
\end{quote}

The verb behind \emph{lleno}, full.

\begin{cousinweb}
\textbf{llenar} is built on \emph{lleno}, \textquotedblleft{}full\textquotedblright{}, from Latin \emph{plenus}, \textquotedblleft{}full\textquotedblright{} — inherited: the \emph{pl-} became \emph{ll-}, as \emph{clamare} became \emph{llamar}.
\end{cousinweb}

\begin{grammarlens}[title={What you've built}]
One more word for reading notices, forms and messages.
\end{grammarlens}

\subsection*{Your turn}
\begin{itemize}
\raggedright
  \item \emph{Say it:} \emph{llenar}
  \item \emph{Say it:} \emph{llenar}, once more
  \item \emph{Say it:} a sentence of your own with \emph{llenar}
  \item \emph{Recall:} say the Spanish for to put, then the Spanish for to take away, then say \emph{llenar} again
\end{itemize}

\begin{tcolorbox}[breakable,colback=teal!4,colframe=teal!35!black,title={Before you move on}]
What does llenar mean? (\textquotedblleft{}To fill\textquotedblright{}.) Say it once more.
\end{tcolorbox}
